\title{Group Sequence Policy Optimization}

\begin{abstract}
In this work, we present Group Sequence Policy Optimization (GSPO), a robust and highly effective reinforcement learning approach tailored for training large language models. Distinct from existing methods that utilize token-wise importance ratios, GSPO instead leverages a sequence-level importance ratio, applying clipping, rewards, and optimization at the granularity of whole sequences. Our empirical results indicate that GSPO offers substantial gains in both training stability and performance when benchmarked against the GRPO algorithm, with significant advances in the stabilization of reinforcement learning involving Mixture-of-Experts (MoE) models, as well as streamlining RL system architectures. These advantages have directly facilitated notable enhancements in the most recent Qwen3 models.
\end{abstract}

\section{Introduction}
\label{sec:intro}

Reinforcement learning (RL) has become a central approach for advancing the capabilities of language models at scale \citep{o1,dpsk-r1,qwq32b,qwen3}. Leveraging large-scale RL enables language models to engage in more extended and intricate reasoning, empowering them to address challenging tasks such as advanced mathematics and programming problems that require multi-step solutions.

For RL scaling to be effective with increased computational resources, ensuring consistent and resilient training behavior is paramount. Yet, leading RL methods like GRPO \citep{grpo} frequently encounter significant stability challenges when applied to very large language models, which may culminate in dramatic and often irreversible degradation of model performance, referred to as model collapse \citep{qwen3, minimax-m1}. Such instability imposes a critical bottleneck for extending the potential of language models via further RL optimization.

In this work, we demonstrate that the core source of instability in GRPO arises from an improper use and subsequent invalidation of importance sampling weights within its algorithmic structure. This flaw injects high-variance noise into the training process, which not only accumulates as the response length increases but is also exacerbated by the clipping strategy, eventually resulting in the collapse of the model.

To overcome these principal shortcomings, we introduce \textbf{Group Sequence Policy Optimization (GSPO)}, a novel reinforcement learning framework for large language model training. GSPO’s primary contribution is a rigorously defined importance ratio, derived from sequence likelihood \citep{click}, thereby restoring adherence to the foundational concept of importance sampling. Furthermore, GSPO calculates normalized rewards by treating them as advantages across several responses for a given query, facilitating coherent optimization at the sequence level with the underlying reward structure.

Through empirical analysis, we find that GSPO substantially outperforms GRPO in terms of training stability, efficiency, and overall effectiveness. Notably, GSPO inherently addresses the stability problems that arise in RL training with large Mixture-of-Experts (MoE) models, removing the reliance on intricate stabilization techniques and pointing towards the possibility of streamlined RL infrastructure. These strengths of GSPO have directly facilitated remarkable enhancements in the performance of the most recent Qwen3 models. We anticipate that GSPO will serve as a reliable and scalable algorithmic basis for ongoing progress in large-scale RL training with language models.

\section{Preliminaries}

\paragraph{Notation}
Throughout this work, we represent an autoregressive language model parameterized by $\theta$ as the policy $\pi_\theta$. The symbol $x$ is reserved for a query, while $\mathcal{D}$ indicates the collection of queries. If $y$ is the response generated for a particular query $x$, the probability of $y$ given $x$ under $\pi_\theta$ is expressed as $\pi_\theta (y | x)=\prod_{t=1}^{|y|} \pi_\theta (y_t | x, y_{<t} )$, where $|y|$ refers to the token length of $y$. The evaluation of a query-response tuple $(x, y)$ is performed via a verifier $r$, with the resulting reward given by $r(x, y) \in [0, 1]$.

\paragraph{Proximal Policy Optimization (PPO)} Using samples generated from the old policy $\pi_{\theta_\text{old}}$, PPO \citep{ppo} constrains the policy update within a proximal region of the old policy through the clipping mechanism. Specifically, PPO employs the following objective for policy optimization (we omit the KL regularization term hereinafter for brevity, as it is not the focus of this paper):

\begin{align}
\mathcal{J}_\text{PPO}(\theta) = \mathbb{E}_{ x \sim \mathcal{D},\, y \sim \pi_{\theta_\text{old}}( \cdot | x) }
\left[ \frac{1}{|y|} \sum_{t=1}^{|y|} 
\min \left( w_{t}(\theta) \widehat{A}_{t},  \, \mathrm{clip} \left( w_{t}(\theta), 1 - {\varepsilon}, 1 + {\varepsilon}\right) \widehat{A}_{t} \right)
\right],
\end{align}

Here, the token importance ratio $w_{t}(\theta)$ is given by $
w_{t}(\theta) = \frac{ \pi_{\theta} (y_{t} | x, y_{<t}) }{ \pi_{\theta_\text{old}} (y_{t} | x,y_{<t})}$, while the advantage $\widehat{A}_{t}$ corresponding to $y_t$ is produced by a separate value model, and $\varepsilon$ denotes the permissible range for clipping these importance weights.

A primary obstacle encountered when applying PPO is its strong dependence on the value model. Typically, this model is of comparable complexity to the policy model, resulting in significant computational and memory requirements. Additionally, the algorithm's overall performance is closely tied to the accuracy of value predictions. Not only is developing a dependable value model intrinsically difficult, but making sure this model scales for longer sequences and intricate tasks exacerbates the issue.

\paragraph{Group Relative Policy Optimization (GRPO)} GRPO \citep{grpo} bypasses the need for the value model by computing the relative advantage of each response within a group of responses to the same query. Specifically, GRPO optimizes the following objective:

\begin{align}
\label{equ:grpo}
\mathcal{J}_\text{GRPO}(\theta) = \mathbb{E}_{ x \sim \mathcal{D},\, \{y_i\}_{i=1}^G \sim \pi_{\theta_\text{old}}( \cdot | x) }
\left[ \frac{1}{G} \sum_{i=1}^{G} \frac{1}{|y_i|} \sum_{t=1}^{|y_i|} 
\min \left( w_{i,t}(\theta) \widehat{A}_{i,t},  \, \mathrm{clip} \left( w_{i,t}(\theta), 1 - {\varepsilon}, 1 + {\varepsilon}\right) \widehat{A}_{i,t} \right)
\right],
\end{align}

where $G$ is the number of generated responses to each query $x$ (i.e., the group size), and the importance ratio $w_{i,t}(\theta)$ and advantage $\widehat{A}_{i,t}$ of token $y_{i,t}$ are:

\begin{align}
    w_{i,t}(\theta)=\frac{ \pi_{\theta} (y_{i,t} | x, y_{i,<t}) }{ \pi_{\theta_\text{old}} (y_{i,t} | x,y_{i,<t})},\quad\ 
    \widehat{A}_{i,t} = \widehat{A}_{i} = \frac{r(x, y_i) - \mathrm{mean} \left( \{ r(x, y_i) \}_{i=1}^G \right) }{ \mathrm{std} \left( \{ r(x, y_i) \}_{i=1}^G \right) },
\end{align}

respectively, where all the tokens in $y_i$ share the same advantage as $\widehat{A}_{i}$.

\section{Motivation}
\label{sec:motivation}

As models become larger, incorporate sparsity (such as in Mixture-of-Experts architectures), and produce longer responses, it is increasingly necessary to employ substantial rollout batch sizes to achieve optimal hardware performance in reinforcement learning (RL). To enhance the efficiency with which data is used, practitioners commonly divide these substantial batches into several smaller mini-batches when computing gradient updates. This practice results in an inherently off-policy scenario, wherein responses $y$ are drawn from a prior policy $\pi_{\theta_\text{old}}$ rather than the updated policy $\pi_\theta$ that is currently being trained. Consequently, this off-policy discrepancy underscores the importance of the clipping strategies used in PPO and GRPO; these mechanisms restrict the impact of samples that diverge significantly from the current policy during gradient calculations.

While mechanisms like clipping aim to manage this off-policy discrepancy, we identify a more fundamental issue in GRPO: \textit{its objective is ill-posed}. This problem becomes particularly acute when training large models on long-response tasks, leading to catastrophic model collapse. The ill-posed nature of the GRPO objective stems from a misapplication of importance sampling weights. The principle of importance sampling is to estimate the expectation of a function $f$ under a target distribution $\pi_\text{tar}$ by re-weighting samples drawn from a behavior distribution $\pi_\text{beh}$:

\begin{align}
\mathbb{E}_{z \sim \pi_\text{tar}} \left[ f(z) \right] = \mathbb{E}_{z \sim \pi_\text{beh}} \left[ \frac{ \pi_\text{tar} (z) }{ \pi_\text{beh} (z)} \, f(z) \right].
\end{align}

A critical aspect of this process is that it depends on aggregating results from a large number of samples ($N \gg 1$) drawn from the behavior distribution $\pi_\text{beh}$, in order for the importance weight $\frac{ \pi_\text{tar} (z) }{ \pi_\text{beh} (z)}$ to reliably account for the mismatch between the distributions.

By contrast, GRPO introduces the importance weight $\frac{ \pi_{\theta} (y_{i,t} | x, y_{i,<t}) }{ \pi_{\theta_\text{old}} (y_{i,t} | x,y_{i,<t})}$ separately at each token step $t$. This weight, computed using only a single instance $y_{i,t}$ sampled from the next-token distribution $\pi_{\theta_\text{old}}( \cdot | x, y_{i, <t})$, does not adequately serve its intended purpose of adjusting for off-policy sampling. Rather, it injects significant stochastic noise into the gradient updates, especially for longer sequences, an effect aggravated by the use of the clipping operation. Through empirical analysis, we have found that such a strategy can frequently cause catastrophic model collapse—an effect that proves persistent. Once collapse has occurred, neither resuming from earlier saved model states nor careful adjustment of hyperparameters (such as changing clipping thresholds, varying the length of generated sequences, or altering the RL reward queries) can restore the model’s performance.

These empirical findings reveal an intrinsic flaw in how GRPO is formulated. Specifically, the ineffectiveness of the token-level importance weight highlights a more general guideline: \textbf{optimization objectives should be defined at the same granularity as the reward signal}. Since the reward mechanism evaluates entire sequences, it is ill-advised to apply off-policy corrections at the level of individual tokens. This insight inspires a shift away from the token-level formulation, motivating us to investigate the application of importance weights and the construction of optimization criteria at the \textit{sequence level} instead.

\section{Algorithm}

\subsection{GSPO: Group Sequence Policy Optimization}

Although the token-level importance weight $\frac{ \pi_{\theta} (y_{i,t} | x, y_{i,<t}) }{ \pi_{\theta_\text{old}} (y_{i,t} | x,y_{i,<t})}$ presents difficulties within GRPO, we find that the \textit{sequence-level} importance weight $\frac{ \pi_{\theta} (y | x) }{ \pi_{\theta_\text{old}} (y | x)}$ acquires a well-defined theoretical interpretation in the setting of language generation. Specifically, it quantifies the extent to which a response $y$—sampled under $\pi_{\theta_\text{old}} (\cdot | x)$—differs from the distribution induced by $\pi_{\theta} (\cdot | x)$, coherently corresponding to the sequence-level reward, and simultaneously providing an informative criterion for applying the clipping strategy.

Based on this straightforward observation, we propose the \textbf{Group Sequence Policy Optimization (GSPO)} algorithm. GSPO employs the following sequence-level optimization objective:

\begin{align}
\mathcal{J}_\text{GSPO} (\theta) =
\mathbb{E}_{ x \sim \mathcal{D},\, \{y_i\}_{i=1}^G \sim \pi_{\theta_\text{old}}( \cdot | x) }
\left[ 
\frac{1}{G} \sum_{i=1}^{G}
\min \left( s_{i}(\theta)  \widehat{A}_{i},  \, \mathrm{clip} \left( s_{i}(\theta), 1 - {\varepsilon}, 1 + {\varepsilon} \right) \widehat{A}_{i} \right) 
\right],
\label{equ:gspo}
\end{align}

where we adopt the group-based advantage estimation:

\begin{align}
\widehat{A}_{i} = \frac{r(x, y_i) - \mathrm{mean} \left( \{ r(x, y_i) \}_{i=1}^G \right) }{ \mathrm{std} \left( \{ r(x, y_i) \}_{i=1}^G \right) },
\end{align}

and define the importance ratio $s_{i}(\theta)$ based on sequence likelihood \citep{click}:

\begin{align}
s_{i}(\theta) = \left( \frac{ \pi_{\theta} (y_i | x) }{ \pi_{\theta_\text{old}} (y_i | x)} \right)^{\frac{1}{|y_i|}}
=
\exp \left( \frac{1}{|y_i|} \sum_{t=1}^{|y_i|} \log \frac{ \pi_{\theta} (y_{i,t} | x, y_{i,<t}) }{ \pi_{\theta_\text{old}} (y_{i,t} | x,y_{i,<t})} \right).
\end{align}

Consequently, GSPO implements clipping at the response level rather than on a per-token basis, thereby removing excessively “off-policy” samples from contributing to the gradient estimation process, which is consistent with the sequence-level approach to both rewarding and optimization. Additionally, we incorporate length normalization in $s_{i}(\theta)$ to mitigate variance and to maintain $s_{i}(\theta)$ within a consistent numerical scale. Without such normalization, changes in the likelihoods of only a handful of tokens could provoke significant instability in the sequence-level importance ratio, necessitating individualized clipping thresholds for responses of varying lengths. It should also be pointed out that, as a result of differing importance ratio definitions, the appropriate clipping ranges for GSPO and earlier methods (such as GRPO) often diverge substantially, typically by several orders of magnitude.

\subsection{Gradient Analysis}
\label{subsec:gradient}

We can derive the gradient of the GSPO objective as follows (clipping is omitted for brevity):

\begin{align}
\nabla_{\theta} \mathcal{J}_\text{GSPO} (\theta)
=&\ 
\nabla_{\theta} \mathbb{E}_{ x \sim \mathcal{D},\, \{y_i\}_{i=1}^G \sim \pi_{\theta_\text{old}}( \cdot | x) }
\left[ 
\frac{1}{G} \sum_{i=1}^{G} s_{i}(\theta) \widehat{A}_{i}
\right] \\
= &\ 
\mathbb{E}_{ x \sim \mathcal{D},\, \{y_i\}_{i=1}^G \sim \pi_{\theta_\text{old}}( \cdot | x) }
\left[ 
\frac{1}{G} \sum_{i=1}^{G}
s_{i}(\theta) \widehat{A}_{i}
\cdot \nabla_{\theta} \log s_{i}(\theta)
\right] \\
=&\ 
\mathbb{E}_{ x \sim \mathcal{D},\, \{y_i\}_{i=1}^G \sim \pi_{\theta_\text{old}}( \cdot | x) }
\left[ 
\frac{1}{G} \sum_{i=1}^{G} 
\left( \frac{ \pi_{\theta} (y_i | x) }{ \pi_{\theta_\text{old}} (y_i | x)} \right)^{\frac{1}{|y_i|}} \widehat{A}_{i} 
\cdot \frac{1}{|y_i|} \sum_{t=1}^{|y_i|} \nabla_{\theta} \log \pi_{\theta} (y_{i,t} | x, y_{i,<t}) 
\right].
\label{equ:gspo_grad}
\end{align}

For comparison, the gradient of the GRPO objective is as follows (note that $\widehat{A}_{i,t} = \widehat{A}_{i}$):

\begin{align}
\nabla_{\theta} \mathcal{J}_\text{GRPO} (\theta) 
=&\ 
\nabla_{\theta} \mathbb{E}_{ x \sim \mathcal{D},\, \{y_i\}_{i=1}^G \sim \pi_{\theta_\text{old}}( \cdot | x) }
\left[ \frac{1}{G} \sum_{i=1}^{G} \frac{1}{|y_i|} \sum_{t=1}^{|y_i|} 
w_{i,t}(\theta) \widehat{A}_{i,t}
\right] \\
=&\
\mathbb{E}_{ x \sim \mathcal{D},\, \{y_i\}_{i=1}^G \sim \pi_{\theta_\text{old}}( \cdot | x) }
\left[ \frac{1}{G} \sum_{i=1}^{G} \widehat{A}_{i} 
\cdot \frac{1}{|y_i|} \sum_{t=1}^{|y_i|} 
\frac{ \pi_{\theta} (y_{i,t} | x, y_{i,<t}) }{ \pi_{\theta_\text{old}} (y_{i,t} | x,y_{i,<t})} 
\nabla_{\theta} \log \pi_{\theta} (y_{i,t} | x, y_{i,<t})  
\right].
\end{align}

The primary difference between GSPO and GRPO is rooted in \textit{the manner in which they assign weights to the gradients of token log likelihoods}. In the case of GRPO, individual tokens are multiplied by an ``importance weight'' defined by $\frac{ \pi_{\theta} (y_{i,t} | x, y_{i,<t}) }{ \pi_{\theta_\text{old}} (y_{i,t} | x,y_{i,<t})}$. These variable weights can range within $(0, 1+\varepsilon]$ when $\widehat{A}_{i} >0$, or $[1-\varepsilon, +\infty)$ when $\widehat{A}_{i} < 0$, and are significant enough that, over time, their cumulative effect can introduce instability and lead to unexpected behaviors during model training. By contrast, GSPO applies uniform weighting across all tokens in a sequence, thereby removing this source of volatility present in GRPO.

\subsection{GSPO-token: A Token-level Objective Variant}

In scenarios like multi-turn RL, we may desire a finer-grained advantage adjustment than the sequence level. To this end, we introduce a token-level objective variant of GSPO, namely \textbf{GSPO-token}, to allow token-wise advantage customization:

\begin{align}
\mathcal{J}_\text{GSPO-token}(\theta)
=
\mathbb{E}_{ x \sim \mathcal{D},\, \{y_i\}_{i=1}^G \sim \pi_{\theta_\text{old}}( \cdot | x) }
\left[ 
\frac{1}{G} \sum_{i=1}^{G} \frac{1}{|y_i|} \sum_{t=1}^{|y_i|} 
\min \left( s_{i,t}(\theta) \widehat{A}_{i,t},  \, \mathrm{clip} \left( s_{i,t}(\theta), 1 - {\varepsilon}, 1 + {\varepsilon} \right) \widehat{A}_{i,t} \right) 
\right],
\label{equ:gspo-token}
\end{align}

where

\begin{align}
s_{i,t}(\theta) 
= \mathrm{sg} \left[ s_{i}(\theta) \right]  \cdot \frac{ \pi_{\theta} (y_{i,t} | x, y_{i,<t}) }{ \mathrm{sg} \left[ \pi_{\theta} (y_{i,t} | x, y_{i,<t}) \right] },
\end{align}

and $\mathrm{sg}[\cdot]$ denotes only taking the numerical value but stopping the gradient, corresponding to the \texttt{detach} operation in PyTorch. The gradient of GSPO-token can be derived as:

\begin{align}
\nabla_{\theta} \mathcal{J}_\text{GSPO-token}(\theta)
=&\ 
\nabla_{\theta} \mathbb{E}_{ x \sim \mathcal{D},\, \{y_i\}_{i=1}^G \sim \pi_{\theta_\text{old}}( \cdot | x) }
\left[ 
\frac{1}{G} \sum_{i=1}^{G}
\frac{1}{|y_i|} \sum_{t=1}^{|y_i|} 
s_{i,t}(\theta) \widehat{A}_{i,t}
\right] \\
=&\ 
\mathbb{E}_{ x \sim \mathcal{D},\, \{y_i\}_{i=1}^G \sim \pi_{\theta_\text{old}}( \cdot | x) }
\left[ 
\frac{1}{G} \sum_{i=1}^{G} s_i (\theta) \cdot \frac{1}{|y_i|} \sum_{t=1}^{|y_i|} 
\widehat{A}_{i,t} \frac{ \nabla_{\theta} \pi_{\theta} (y_{i,t} | x, y_{i,<t}) }{ \pi_{\theta} (y_{i,t} | x, y_{i,<t}) }
\right] \\
=&\ 
\mathbb{E}_{ x \sim \mathcal{D},\, \{y_i\}_{i=1}^G \sim \pi_{\theta_\text{old}}( \cdot | x) }
\left[ 
\frac{1}{G} \sum_{i=1}^{G}  \left( \frac{ \pi_{\theta} (y_i | x) }{ \pi_{\theta_\text{old}} (y_i | x)} \right)^{\frac{1}{|y_i|}}
\cdot \frac{1}{|y_i|} \sum_{t=1}^{|y_i|}
\widehat{A}_{i,t} \nabla_{\theta} \log \pi_{\theta} (y_{i,t} | x, y_{i,<t})  
\right].
\label{equ:gspo-token_grad}
\end{align}

Observe that the factor $\frac{ \pi_{\theta} (y_{i,t} | x, y_{i,<t}) }{ \mathrm{sg} \left[ \pi_{\theta} (y_{i,t} | x, y_{i,<t}) \right] }$ evaluates numerically to 1, implying that $s_{i,t}(\theta)$ and $s_{i}(\theta)$ yield equivalent values. Upon inspection of Equation~\eqref{equ:gspo} alongside~\eqref{equ:gspo-token}, as well as Equation~\eqref{equ:gspo_grad} compared to~\eqref{equ:gspo-token_grad}, it becomes clear that GSPO-token and GSPO coincide exactly in terms of their optimization objectives, clipping criteria, and theoretical gradients, provided all tokens in the sequence $y_i$ share identical advantage values (that is, $\widehat{A}_{i,t} = \widehat{A}_i$ for all $t$). Nevertheless, GSPO-token distinguishes itself by allowing independent adjustment of advantages at the token level, offering greater granularity and flexibility.

\section{Experiments and Discussion}

\subsection{Empirical Results}

We conduct experiments utilizing a cold-start model, initialized from Qwen3-30B-A3B-Base, and present both the reward trajectories during training and the performance trends on the AIME'24 (mean Pass@1 across 32 samples), LiveCodeBench (202410-202502, mean Pass@1 across 8 samples), and CodeForces (Elo Rating) evaluation suites. For RL training, rollout batches are divided into four mini-batches to perform gradient updates. With GSPO, we specify the clipping bounds in Equation~\eqref{equ:gspo} as 3e-4 (left) and 4e-4 (right). Our baseline for comparison is GRPO, where the clipping thresholds in Equation~\eqref{equ:grpo} are fixed at 0.2 (left) and 0.27 (right)—parameters optimized for an equitable comparison. Importantly, GRPO relies on the Routing Replay training approach to ensure proper convergence in MoE RL, a topic we will examine further in \S~\ref{subsec:routing_replay}. In contrast, \textit{GSPO eliminates the reliance on this strategy}.

Figure~\ref{fig:results} illustrates that training with GSPO exhibits steady progress across the training process. Notably, we find that \textbf{GSPO consistently enhances performance as more training compute is allocated, the query set is periodically refreshed, and the sequence generation length is increased}. Additionally, GSPO outperforms GRPO in terms of training efficiency, achieving superior accuracy and benchmark results with equivalent training compute and query consumption. Lastly, GSPO has been effectively utilized in RL training of the latest Qwen3 models, providing strong evidence for its capacity to maximize the scaling benefits of RL in large language models.

\subsection{Curious Observation on Clipping Fractions}

One principal difference between GSPO and GRPO lies in the clipping methodology: GSPO applies clipping at the response level, whereas GRPO performs clipping on a per-token basis. As illustrated in Figure~\ref{fig:clipping}, this approach leads to GSPO clipping token fractions that are two orders of magnitude greater than those in GRPO, regardless of adjustments to the clipping intervals, which do not diminish this substantial gap. Notably, although GSPO discards many more tokens—thereby relying on a reduced set for model updates or gradient calculations—it nonetheless surpasses GRPO in terms of training efficiency. This seemingly paradoxical outcome, where training becomes more effective with the exclusion of a larger portion of the data, suggests that the gradient estimates derived by GRPO at the token level are intrinsically noisy and less capable of leveraging sampled data. Conversely, GSPO’s sequence-based method delivers a stronger and more consistent learning signal.

\subsection{Benefit of GSPO for MoE Training}
\label{subsec:routing_replay}

\paragraph{Background}
While training dense models with RL generally proceeds smoothly, MoE models present distinct challenges due to their inherently sparse activation patterns. Notably, our experiments indicate that applying the GRPO algorithm to MoE architectures results in significant instability related to \textit{expert-activation volatility}, which undermines RL training convergence. Specifically, we observed that expert selections for identical outputs can vary widely after one or more gradient updates. As an illustration, for the 48-layer Qwen3-30B-A3B-Base architecture, approximately 10\% of the experts triggered for a particular rollout sample change following each RL gradient update, when comparing selections made by the updated policy $\pi_{\theta}$ against those of the prior policy $\pi_{\theta_\text{old}}$. This issue intensifies in larger, deeper MoE networks, causing the token-wise importance ratios $w_{i,t}(\theta) = \frac{ \pi_{\theta} (y_{i,t} | x, y_{i,<t}) }{ \pi_{\theta_\text{old}} (y_{i,t} | x,y_{i,<t})}$ to display severe volatility, thus invalidating these metrics as outlined in \S~\ref{sec:motivation} and~\ref{subsec:gradient}. As a result, RL training convergence is significantly impeded.

\paragraph{Our Previous Approach} In our earlier work addressing this issue, we adopted the \textbf{Routing Replay} training method. In particular, we store the expert routing decisions made by $\pi_{\theta_\text{old}}$ and ``replay'' these same routes using $\pi_{\theta}$ while evaluating the importance ratios $w_{i,t}(\theta) = \frac{ \pi_{\theta} (y_{i,t} | x, y_{i,<t}) }{ \pi_{\theta_\text{old}} (y_{i,t} | x,y_{i,<t})}$. This strategy guarantees that both $\pi_{\theta} (y_{i,t} | x, y_{i,<t})$ and $\pi_{\theta_\text{old}} (y_{i,t} | x,y_{i,<t})$ utilize the identical network structure for each token $y_{i,t}$, thereby preserving the stability of the importance ratio at the token level and ensuring that training consistently updates the same network path. As shown in Figure~\ref{fig:routing_replay}, Routing Replay proves to be a vital component for achieving stable GRPO training and facilitating the effective convergence of MoE models.

\paragraph{Benefit of GSPO} While Routing Replay allows for the stable GRPO training of MoE models, its reliance on reusing routing modes introduces added memory and communication costs and can restrict the model's effective capacity. GSPO, as illustrated in Figure~\ref{fig:results}, circumvents these limitations by dispensing with Routing Replay altogether. It is fully compatible with conventional computation of the importance ratios $s_i(\theta)$, achieves standard convergence, and maintains stable optimization. Crucially, GSPO’s approach centers exclusively on the sequence likelihood (i.e., $\pi_{\theta} (y_i | x)$), showing little sensitivity to the likelihood of individual tokens (i.e., $\pi_{\theta} (y_{i,t} | x, y_{i,<t})$). Because the MoE model consistently preserves its language modeling performance, the overall sequence likelihood remains steady. Ultimately, GSPO addresses the fluctuations in expert activation in MoE models at a fundamental level, rendering intricate strategies like Routing Replay unnecessary. This not only streamlines and stabilizes training, but also empowers the model to utilize its full potential unimpeded by artificial restrictions.

\subsection{Benefit of GSPO for RL Infrastructure}

Due to differences in numerical precision between training engines (such as Megatron) and inference engines (such as SGLang and vLLM), it is common practice to recompute the likelihoods of sampled responses from the old policy $\pi_{\theta_\text{old}}$ using the training engine. In contrast, GSPO relies on sequence-level likelihoods for optimization, which are, by nature, less sensitive to discrepancies in precision compared to token-level likelihoods. As a result, GSPO enables the direct utilization of likelihoods output by the inference engine during optimization, eliminating the necessity to perform additional recomputation on the training engine. This property is particularly advantageous in contexts like partial rollout, multi-turn RL, and training-inference disaggregation frameworks.

\section{Conclusion}

We introduce Group Sequence Policy Optimization (GSPO), a novel reinforcement learning method tailored for large language model training. Built upon the central concept of importance sampling, GSPO establishes importance ratios derived from sequence likelihood and implements clipping, rewarding, and optimization at the sequence level. Compared to GRPO, GSPO offers significantly enhanced training stability, improved efficiency, and better overall performance, particularly excelling in large-scale reinforcement learning for MoE models. This approach serves as the underlying driver for the remarkable advancements observed in recent Qwen3 models. With GSPO as a robust and scalable foundation, we aim to further expand reinforcement learning and anticipate substantial breakthroughs in the development of intelligence.

\end{document}